\documentclass[a4paper,12pt]{article}

\usepackage[T1]{fontenc}
\usepackage[utf8]{inputenc}
\usepackage[polish]{babel}
\usepackage{lmodern}
\usepackage{textcomp}
\usepackage{microtype}

\usepackage[a4paper,left=2.5cm,right=2.5cm,top=2.5cm,bottom=2.5cm]{geometry}
\usepackage{setspace}
\onehalfspacing
\usepackage{parskip}

\usepackage{amsmath,amssymb,mathtools}
\usepackage{siunitx}
\sisetup{locale = DE}

\usepackage{graphicx}
\usepackage{float}
\usepackage{caption}
\usepackage{subcaption}
\usepackage{booktabs}
\usepackage{longtable}
\usepackage{array}
\usepackage{tabularx}
\usepackage{multirow}
\usepackage[table]{xcolor}

\usepackage{enumitem}
\usepackage[hidelinks]{hyperref}
\usepackage[nameinlink,capitalise,noabbrev]{cleveref}

\usepackage{fancyhdr}
\pagestyle{fancy}
\fancyhf{}
\lhead{Interfejsy cyfrowe}
\rhead{Projekt}
\cfoot{\thepage}
\renewcommand{\headrulewidth}{0.4pt}
\setlength{\headheight}{15pt}

\usepackage{titlesec}
\titleformat{\section}{\Large\bfseries}{\thesection.}{0.7em}{}
\titleformat{\subsection}{\large\bfseries}{\thesubsection.}{0.7em}{}

\usepackage{listings}
\usepackage{courier}

\definecolor{codebg}{RGB}{248,248,248}
\definecolor{codeframe}{RGB}{220,220,220}
\definecolor{codekw}{RGB}{0,70,140}
\definecolor{codecom}{RGB}{0,120,70}
\definecolor{codestr}{RGB}{160,50,20}

\lstdefinestyle{CStyle}{
    language=C,
    inputencoding=utf8,
    extendedchars=true,
    basicstyle=\footnotesize\ttfamily,
    keywordstyle=\color{codekw}\bfseries,
    commentstyle=\color{codecom}\itshape,
    stringstyle=\color{codestr},
    backgroundcolor=\color{codebg},
    frame=single,
    rulecolor=\color{codeframe},
    numbers=left,
    numberstyle=\scriptsize,
    stepnumber=1,
    numbersep=8pt,
    tabsize=4,
    showstringspaces=false,
    breaklines=true,
    captionpos=b
}

\title{\textbf{Projekt interaktywnego systemu sterowania\\z wyświetlaczem LCD, joystickiem, Bluetooth i serwem}}
\author{Jarosław Rybak \\ nr albumu: 275773}
\date{Wrocław, 2026}

\begin{document}

\begin{titlepage}
    \centering
    \vspace*{1cm}

    {\Large \textbf{Politechnika Wrocławska}}\\[0.2cm]
    {\large \textbf{Wydział Elektroniki, Fotoniki i Mikrosystemów}}\\[2.5cm]

    {\huge \textbf{Interfejsy cyfrowe}}\\[0.5cm]
    {\LARGE Sprawozdanie projektowe}\\[1.5cm]

    {\Large \textbf{Interaktywny sterownik mikrokontrolerowy\\z LCD, joystickiem i komunikacją Bluetooth}}\\[3cm]

    \begin{flushright}
        \begin{tabular}{rl}
            Autor: & Jarosław Rybak \\
            Nr albumu: & 275773 \\
            Kierunek: & Elektroniczne Systemy Mechatroniki \\
            Prowadzący: & dr inż. Mirosław Gierczak
        \end{tabular}
    \end{flushright}

    \vfill
    {\large Wrocław 2026}
\end{titlepage}

\tableofcontents
\newpage

\section{Wprowadzenie}

Celem projektu było uruchomienie interaktywnego systemu mikrokontrolerowego opartego na układzie AVR, który łączy kilka typowych interfejsów cyfrowych i analogowych. W jednej aplikacji połączono wyświetlacz LCD 2x16, joystick analogowy, komunikację Bluetooth oraz sterowanie serwomechanizmem.

Projekt ma charakter demonstracyjny i laboratoryjny. Układ można obsługiwać lokalnie przez joystick lub zdalnie przez telefon z aplikacją terminala Bluetooth. System reaguje na komendy tekstowe, wyświetla komunikaty na LCD oraz odsyła dane pomiarowe przez UART.

\section{Cel projektu}

Głównym celem projektu było praktyczne przećwiczenie:

\begin{itemize}[leftmargin=1.2cm]
    \item obsługi portów wejścia i wyjścia mikrokontrolera,
    \item sterowania wyświetlaczem LCD w trybie 4-bitowym,
    \item wykorzystania przetwornika ADC do odczytu joysticka,
    \item konfiguracji sprzętowego UART i współpracy z modułem Bluetooth,
    \item generacji sygnału PWM do sterowania serwem,
    \item organizacji programu w postaci prostego menu i osobnych trybów pracy.
\end{itemize}

\section{Opis funkcjonalny systemu}

Po uruchomieniu programu wykonywana jest inicjalizacja LCD, ADC, UART oraz Timer1 dla serwomechanizmu. Następnie użytkownik trafia do menu głównego sterowanego ruchem joysticka w osi Y. Wejście do wybranego trybu następuje po naciśnięciu przycisku joysticka.

Aktualna wersja programu zawiera następujące tryby:

\begin{itemize}[leftmargin=1.2cm]
    \item \textbf{O mnie} --- prezentacja danych autora,
    \item \textbf{Aplikacja} --- ekran informacyjny projektu,
    \item \textbf{BT TEXT} --- wyświetlanie tekstu wpisanego z telefonu,
    \item \textbf{BT CTRL} --- obsługa komend sterujących i diagnostycznych,
    \item \textbf{JOY MONITOR} --- podgląd wartości osi X i Y oraz kierunku wychylenia,
    \item \textbf{SERVO TEST} --- sterowanie serwem za pomocą joysticka.
\end{itemize}

\section{Wykorzystane elementy}

W projekcie zastosowano następujące elementy:

\begin{itemize}[leftmargin=1.2cm]
    \item mikrokontroler ATmega32A,
    \item zestaw uruchomieniowy ATB 1.05,
    \item wyświetlacz LCD 2x16 zgodny z HD44780,
    \item joystick analogowy,
    \item moduł Bluetooth BTM222B,
    \item serwomechanizm sterowany sygnałem PWM,
    \item przewody połączeniowe oraz zasilanie pomocnicze.
\end{itemize}

\section{Połączenia sprzętowe}

\subsection{LCD}

\begin{itemize}[leftmargin=1.2cm]
    \item RS $\rightarrow$ PB0,
    \item E $\rightarrow$ PB1,
    \item D4 $\rightarrow$ PB2,
    \item D5 $\rightarrow$ PB3,
    \item D6 $\rightarrow$ PB4,
    \item D7 $\rightarrow$ PB5,
    \item RW $\rightarrow$ GND.
\end{itemize}

\subsection{Joystick}

\begin{itemize}[leftmargin=1.2cm]
    \item VRx $\rightarrow$ PA2 / ADC2,
    \item VRy $\rightarrow$ PA3 / ADC3,
    \item SW $\rightarrow$ PD5,
    \item zasilanie $\rightarrow$ +5V oraz GND.
\end{itemize}

\subsection{Bluetooth}

\begin{itemize}[leftmargin=1.2cm]
    \item TX modułu $\rightarrow$ PD0 / RX,
    \item RX modułu $\rightarrow$ PD1 / TX,
    \item VCC $\rightarrow$ +5V,
    \item GND $\rightarrow$ GND.
\end{itemize}

\subsection{Serwo}

\begin{itemize}[leftmargin=1.2cm]
    \item sygnał sterujący $\rightarrow$ PD4 / OC1B,
    \item masa serwa $\rightarrow$ wspólna masa układu,
    \item zasilanie serwa $\rightarrow$ zgodnie z wymaganiami wykonawczymi stanowiska.
\end{itemize}

\section{Komunikacja Bluetooth}

Tryb BT CTRL pozwala sterować systemem z poziomu telefonu. W projekcie zaimplementowano następujące komendy:

\begin{lstlisting}[style=CStyle,caption={Komendy obslugiwane przez aplikacje}]
1      -> pomoc
2tekst -> wyswietl tekst na LCD
3      -> jednorazowy raport joysticka
4      -> wlacz stream joysticka
5      -> wylacz stream joysticka
6      -> pokaz dane autora
7      -> wyczysc LCD
81..89 -> ustaw poziom serwa 1..9
90     -> ustaw serwo w pozycji centralnej
\end{lstlisting}

Przez UART wysyłane są odpowiedzi statusowe, potwierdzenia wykonania oraz raporty zawierające wartości osi joysticka i stan przycisku.

\section{Struktura programu}

Kod źródłowy został uporządkowany w trzech podstawowych plikach:

\begin{itemize}[leftmargin=1.2cm]
    \item \texttt{main.c} --- logika programu, menu, ADC, UART, Bluetooth i serwo,
    \item \texttt{lcd.c} --- implementacja sterownika LCD,
    \item \texttt{lcd.h} --- deklaracje funkcji modułu LCD.
\end{itemize}

W kodzie można wyróżnić sekcje odpowiedzialne za inicjalizację, pętlę główną, logikę trybów, obsługę UART, obsługę ADC oraz generację PWM dla serwomechanizmu.

\section{Uwagi implementacyjne}

W końcowej wersji projektu przyjęto mapowanie linii LCD na port B. Wcześniejsze warianty robocze wykorzystywały inne podłączenie, ale do sprawozdania i repozytorium została przeniesiona wersja odpowiadająca finalnie działającemu stanowisku.

Do programu dodano również osobny tryb testowy serwa oraz dwa niezależne tryby komunikacji Bluetooth: prosty odbiór tekstu oraz tryb komend sterujących.

\section{Wnioski}

Projekt potwierdził możliwość połączenia w jednym układzie kilku typowych interfejsów spotykanych w systemach wbudowanych. W praktyce udało się zintegrować lokalne sterowanie joystickiem, komunikację z telefonem, wyświetlanie informacji na LCD oraz sterowanie serwomechanizmem.

Układ może stanowić bazę do dalszej rozbudowy, na przykład o dodatkowe ekrany menu, nowe komendy Bluetooth, rozbudowane sterowanie wyjściami lub własną aplikację mobilną.

\section*{Miejsce na uzupelnienia przed oddaniem}

Przed finalnym oddaniem sprawozdania warto dodać:

\begin{itemize}[leftmargin=1.2cm]
    \item zdjęcia rzeczywistego układu,
    \item zrzuty ekranu z terminala Bluetooth,
    \item schemat blokowy lub schemat połączeń,
    \item wybrane listingi kodu źródłowego.
\end{itemize}

\end{document}